%% ========================================================================
%% chapter04-directory-structure.tex
%% Struktur Direktori dan Operasi File
%% ========================================================================

\chapter{Struktur Direktori dan Operasi File}
\label{ch:filesystem-ops}

\chapterhero{orange}{Jelajahi struktur sistem berkas Linux, pahami jenis file dan tautan, lalu kuasai pencarian serta pengarsipan untuk pekerjaan harian.}{FHS}{tautan}{arsip}

Bayangkan kamu baru pindah ke rumah kontrakan baru. Ada dapur untuk menyimpan bahan makanan,
kamar tidur untuk barang pribadi, gudang untuk alat-alat, dan teras untuk barang sementara
yang akan segera dipindah. Setiap ruangan punya fungsi khusus dan semua orang di rumah itu
tahu ke mana harus pergi ketika butuh sesuatu. Linux bekerja dengan prinsip yang sama.

Sistem file Linux bukan tumpukan berkas yang diletakkan sembarangan. Setiap direktori punya
peran yang sudah ditetapkan. File konfigurasi aplikasi tinggal di \dir{/etc/}, catatan
aktivitas sistem tersimpan di \dir{/var/log/}, dan program yang bisa dijalankan oleh semua
pengguna ada di \dir{/bin/}. Ketika kamu tahu aturan ini, mencari file di Linux menjadi
pekerjaan yang cepat dan sistematis, bukan tebak-tebakan.

Selain tata letak direktori, Linux membedakan berkas berdasarkan jenisnya: berkas biasa,
direktori, tautan simbolis, berkas perangkat, dan sebagainya. Perbedaan ini bukan sekadar
label, melainkan menentukan bagaimana sistem memperlakukan setiap objek. Memahami perbedaan
ini membuat kamu mampu mendiagnosis masalah akses, menyusun struktur folder proyek yang
rapi, dan bekerja dengan sistem Linux secara lebih percaya diri.

\textbf{Yang Akan Kamu Pelajari}

Setelah menyelesaikan bab ini, kamu akan mampu:

\begin{enumerate}
    \item menjelaskan fungsi direktori-direktori utama dalam hierarki sistem file Linux;
    \item membedakan jenis-jenis berkas di Linux berdasarkan karakter pertama output \cmd{ls -l};
    \item menjelaskan perbedaan hard link dan symbolic link serta membuatnya dengan perintah yang tepat;
    \item mencari berkas berdasarkan nama, tipe, ukuran, dan waktu modifikasi menggunakan \cmd{find} dan \cmd{locate};
    \item membuat, memeriksa, dan mengekstrak arsip terkompresi dengan \cmd{tar} dan \cmd{gzip}.
\end{enumerate}

\begin{notebox}
Seluruh praktek pada bab ini dijalankan pada direktori \dir{~/lab-os/chapter04-files/}.
Buat direktori ini sebelum memulai:
\begin{lstlisting}[language=bash]
mkdir -p ~/lab-os/chapter04-files
cd ~/lab-os/chapter04-files
\end{lstlisting}
\end{notebox}

% ------------------------------------------------------------------------
\section{Standar Hierarki Sistem File}
\label{sec:fhs}

Linux mengatur semua berkas dalam satu pohon direktori tunggal yang dimulai dari akar,
yaitu \dir{/}. Tidak ada konsep ``drive C'' atau ``drive D'' seperti di Windows. Semua
perangkat penyimpanan, partisi, bahkan filesystem jaringan, dipasang (\textit{mount}) ke
dalam satu pohon yang sama. Standar yang mengatur tata letak ini disebut
\textit{Filesystem Hierarchy Standard} (FHS).

Prinsip utama Unix yang perlu kamu pegang adalah: \textit{everything is a file}. Bukan
hanya dokumen teks yang diperlakukan sebagai berkas. Perangkat keras seperti hard disk
direpresentasikan sebagai berkas di \dir{/dev/}. Informasi tentang proses yang sedang
berjalan bisa dibaca seperti membaca berkas di \dir{/proc/}. Memahami prinsip ini membuka
cara pandang baru tentang bagaimana Linux bekerja.

Berikut direktori-direktori utama yang perlu kamu kenal. Setiap administrator Linux
menghafalkan ini secara alami seiring penggunaan sehari-hari.

\begin{itemize}
    \item \dir{} : titik awal seluruh hierarki. Semua direktori lain berasal dari sini.
    \item \dir{/bin/} : program esensial yang dibutuhkan semua pengguna, seperti \cmd{ls}, \cmd{cp}, \cmd{mv}, dan \cmd{bash}.
    \item \dir{/sbin/} : program administrasi sistem, seperti \cmd{fdisk} dan \cmd{iptables}. Umumnya memerlukan \cmd{sudo}.
    \item \dir{/etc/} : semua file konfigurasi sistem. Berformat teks sehingga bisa dibaca dengan editor biasa.
    \item \dir{/home/} : direktori pribadi setiap pengguna. Pengguna \texttt{alice} memiliki \dir{/home/alice/}.
    \item \dir{/root/} : direktori pribadi pengguna \texttt{root}. Dipisah dari \dir{/home/} untuk keamanan.
    \item \dir{/var/} : data yang terus berubah selama sistem berjalan. Subdirektori paling penting adalah \dir{/var/log/} yang menyimpan log sistem.
    \item \dir{/tmp/} : berkas sementara. Biasanya dikosongkan otomatis saat sistem dinyalakan ulang.
    \item \dir{/usr/} : program dan library untuk pengguna. \dir{/usr/bin/} berisi sebagian besar aplikasi yang kamu instal.
    \item \dir{/opt/} : aplikasi pihak ketiga yang dipasang secara mandiri di luar manajer paket.
    \item \dir{/dev/} : berkas perangkat. \file{/dev/sda} adalah hard disk pertama, \file{/dev/tty} adalah terminal.
    \item \dir{/proc/} : filesystem virtual dari kernel. Membacanya memberikan informasi tentang proses dan konfigurasi kernel saat ini.
    \item \dir{/sys/} : filesystem virtual lain dari kernel, berisi informasi perangkat keras dan driver.
    \item \dir{/mnt/} dan \dir{/media/} : titik pemasangan untuk perangkat eksternal seperti USB atau partisi tambahan.
\end{itemize}

\subsection*{Praktek 4.1 Jelajahi Direktori-Direktori Utama}

Tujuan praktek ini adalah membiasakan mata dengan isi direktori standar Linux dan
menemukan file-file konfigurasi penting secara langsung.

\begin{enumerate}
    \item Tampilkan isi direktori root dan perhatikan nama-nama yang sudah kamu kenal:
\begin{lstlisting}[language=bash, caption={Melihat isi root filesystem}]
ls -lah /
\end{lstlisting}
    \textit{Amati:} Perhatikan kolom pertama setiap baris. Karakter \texttt{d} menandakan direktori, sedangkan \texttt{l} menandakan tautan simbolis.

    \item Lihat isi tiga direktori paling sering dikunjungi administrator:
\begin{lstlisting}[language=bash, caption={Menginspeksi /etc, /var/log, dan /home}]
ls /etc | head -20
ls /var/log | head -20
ls /home
\end{lstlisting}

    \item Baca informasi CPU langsung dari filesystem virtual kernel:
\begin{lstlisting}[language=bash, caption={Membaca informasi sistem dari /proc}]
cat /proc/cpuinfo | head -20
cat /proc/meminfo | head -10
\end{lstlisting}
    \textit{Amati:} File di \dir{/proc/} terasa seperti berkas biasa, padahal isinya dibangkitkan oleh kernel secara langsung, bukan disimpan di disk.

    \item Identifikasi tiga berkas konfigurasi penting di \dir{/etc/}:
\begin{lstlisting}[language=bash, caption={Mencari file konfigurasi penting dengan grep}]
ls /etc | grep -E "^(hostname|hosts|passwd|group|fstab)$"
cat /etc/hostname
\end{lstlisting}
\end{enumerate}

\begin{challengebox}
Gunakan kombinasi \cmd{ls} dan \cmd{grep} (yang sudah kamu pelajari di Bab~2) untuk menemukan
tiga berkas konfigurasi di \dir{/etc/} yang namanya mengandung kata \texttt{ssh}. Setelah
menemukannya, tampilkan 10 baris pertama dari masing-masing berkas menggunakan \cmd{head}.
Dokumentasikan perintah yang kamu gunakan dan tuliskan fungsi singkat setiap berkas tersebut.
\end{challengebox}

% ------------------------------------------------------------------------
\section{Jenis Berkas di Linux}
\label{sec:filetypes}

Ketika kamu menjalankan \cmd{ls -l}, karakter pertama pada setiap baris bukan hiasan.
Karakter itu memberitahu jenis berkas tersebut. Linux mengenal tujuh jenis berkas, dan
penting untuk membedakannya karena perilaku sistem terhadap masing-masing jenis berbeda.

Berikut karakter penanda jenis berkas yang akan sering kamu temui:

\begin{itemize}
    \item \texttt{-} : berkas reguler. Dokumen teks, gambar, binary program, semua masuk kategori ini.
    \item \texttt{d} : direktori. Sebuah kontainer yang bisa berisi berkas dan direktori lain.
    \item \texttt{l} : tautan simbolis (\textit{symbolic link}). Sebuah pintasan yang menunjuk ke berkas atau direktori lain.
    \item \texttt{b} : berkas perangkat blok (\textit{block device}). Contohnya adalah partisi hard disk seperti \file{/dev/sda1}.
    \item \texttt{c} : berkas perangkat karakter (\textit{character device}). Contohnya adalah terminal \file{/dev/tty}.
    \item \texttt{p} : pipa bernama (\textit{named pipe}). Digunakan untuk komunikasi antarproses.
    \item \texttt{s} : soket (\textit{socket}). Digunakan untuk komunikasi jaringan atau antarproses lokal.
\end{itemize}

Selain \cmd{ls -l}, dua perintah lain sangat berguna untuk menginspeksi berkas. Perintah
\cmd{stat} menampilkan metadata lengkap sebuah berkas: ukuran, jumlah hard link, inode,
waktu akses terakhir, waktu modifikasi, dan permission dalam notasi oktal maupun simbolis.
Perintah \cmd{file} menebak tipe konten sebuah berkas berdasarkan isinya, bukan ekstensinya.
Ini berguna ketika kamu menemukan berkas tanpa ekstensi dan ingin tahu apakah itu teks,
binary, atau sesuatu yang lain.

\subsection*{Praktek 4.2 Identifikasi Berbagai Jenis Berkas}

Tujuan praktek ini adalah membedakan jenis berkas berdasarkan output \cmd{ls -l} dan
memahami cara membaca metadata berkas secara menyeluruh.

\begin{enumerate}
    \item Buat beberapa berkas dengan jenis berbeda di direktori kerja:
\begin{lstlisting}[language=bash, caption={Membuat berkas reguler dan direktori}]
cd ~/lab-os/chapter04-files
mkdir -p latihan-jenis
echo "isi berkas teks" > latihan-jenis/catatan.txt
mkdir latihan-jenis/subfolder
ln -s latihan-jenis/catatan.txt latihan-jenis/shortcut.txt
ls -lah latihan-jenis/
\end{lstlisting}
    \textit{Amati:} Perhatikan karakter pertama pada setiap baris. \texttt{-} untuk \file{catatan.txt}, \texttt{d} untuk \file{subfolder}, dan \texttt{l} untuk \file{shortcut.txt}. Pada baris tautan simbolis, kamu juga bisa melihat ke mana ia menunjuk.

    \item Gunakan \cmd{stat} untuk melihat metadata lengkap berkas reguler:
\begin{lstlisting}[language=bash, caption={Melihat metadata lengkap dengan stat}]
stat latihan-jenis/catatan.txt
\end{lstlisting}

    \item Gunakan \cmd{file} untuk mengidentifikasi tipe konten berdasarkan isi berkas:
\begin{lstlisting}[language=bash, caption={Mendeteksi tipe berkas dengan file}]
file latihan-jenis/catatan.txt
file latihan-jenis/subfolder
file latihan-jenis/shortcut.txt
file /bin/ls
\end{lstlisting}
    \textit{Amati:} Perhatikan bagaimana \cmd{file} mendeskripsikan \file{/bin/ls} sebagai \textit{ELF executable}, bukan sekadar ``berkas''. Ini berguna ketika kamu menemukan berkas tanpa ekstensi.

    \item Lihat berkas perangkat blok dan karakter di \dir{/dev/}:
\begin{lstlisting}[language=bash, caption={Melihat berkas perangkat di /dev}]
ls -lah /dev/tty /dev/null /dev/zero
ls -lah /dev/sd* 2>/dev/null | head -5
\end{lstlisting}
\end{enumerate}

\begin{challengebox}
Buat sebuah tautan simbolis bernama \file{home-link} di direktori \dir{~/lab-os/chapter04-files/}
yang menunjuk ke direktori home kamu (\dir{~/}). Verifikasi bahwa tautan tersebut bekerja
dengan menjalankan \cmd{ls -la home-link/} dan membandingkan hasilnya dengan \cmd{ls -la ~/}.
Kemudian gunakan \cmd{stat} untuk memeriksa apakah inode \file{home-link} dan \dir{~/}
sama atau berbeda. Jelaskan mengapa hasilnya demikian.
\end{challengebox}

% ------------------------------------------------------------------------
\section{Hard Links dan Symbolic Links}
\label{sec:links}

Di Linux, satu data bisa memiliki lebih dari satu nama. Inilah ide dasar \textit{link}.
Bayangkan sebuah buku di perpustakaan. Kamu bisa menaruh satu kartu indeks yang langsung
menunjuk ke buku itu, atau membuat salinan entri katalog di tempat lain yang tetap
mengacu ke buku fisik yang sama. Dua pendekatan ini sepadan dengan \textit{hard link}
dan \textit{symbolic link}.

\textbf{Hard link} adalah nama tambahan yang menunjuk ke inode yang sama dengan berkas
aslinya. Selama masih ada satu hard link yang tersisa, data berkas tetap ada di disk.
Karena berbagi inode, perubahan isi melalui salah satu nama langsung terlihat pada nama
lainnya. Hard link umumnya hanya dapat dibuat untuk berkas reguler dan tidak melintasi
filesystem yang berbeda.

\textbf{Symbolic link} atau \textit{symlink} berbeda. Ia bukan nama tambahan untuk inode
yang sama, melainkan berkas kecil yang berisi path tujuan. Karena itu, symlink bisa
menunjuk ke direktori, bisa melintasi filesystem, dan bisa menjadi \textit{broken link}
jika targetnya dihapus atau dipindahkan.

Perintah \cmd{ln sumber target} membuat hard link, sedangkan \cmd{ln -s sumber target}
membuat symbolic link. Gunakan \cmd{ls -li} untuk melihat nomor inode dan memverifikasi
apakah dua nama berbagi inode yang sama.

\subsection*{Praktek 4.3 Bandingkan Hard Link dan Symbolic Link}

Tujuan praktek ini adalah memahami perbedaan hard link dan symbolic link melalui inode,
perubahan isi berkas, dan perilaku saat target asli dihapus.

\begin{enumerate}
    \item Siapkan satu berkas sumber dan buat dua jenis link:
\begin{lstlisting}[language=bash, caption={Membuat hard link dan symbolic link}]
cd ~/lab-os/chapter04-files
mkdir -p latihan-link
echo "versi awal" > latihan-link/catatan.txt
ln latihan-link/catatan.txt latihan-link/catatan-hard.txt
ln -s latihan-link/catatan.txt latihan-link/catatan-soft.txt
ls -li latihan-link/
\end{lstlisting}
    \textit{Amati:} \file{catatan.txt} dan \file{catatan-hard.txt} harus memiliki nomor inode yang sama. Symlink memiliki inode berbeda dan menampilkan panah menuju targetnya.

    \item Ubah isi berkas melalui hard link lalu baca kembali dari nama asli:
\begin{lstlisting}[language=bash, caption={Mengubah isi melalui hard link}]
echo "ditambah lewat hard link" >> latihan-link/catatan-hard.txt
cat latihan-link/catatan.txt
\end{lstlisting}

    \item Ubah target asli dan baca isi melalui symlink:
\begin{lstlisting}[language=bash, caption={Membaca target melalui symbolic link}]
echo "ditambah lewat target asli" >> latihan-link/catatan.txt
cat latihan-link/catatan-soft.txt
\end{lstlisting}

    \item Hapus berkas asli lalu amati perilaku kedua link:
\begin{lstlisting}[language=bash, caption={Menghapus target asli}]
rm latihan-link/catatan.txt
ls -li latihan-link/
cat latihan-link/catatan-hard.txt
cat latihan-link/catatan-soft.txt
\end{lstlisting}
    \textit{Amati:} Hard link masih dapat dibaca karena masih menunjuk ke inode yang sama. Symlink gagal karena path targetnya tidak lagi ada.
\end{enumerate}

\begin{challengebox}
Buat satu berkas bernama \file{laporan.txt}, lalu buat hard link dan symbolic link untuk
berkas itu. Gunakan \cmd{ls -li} untuk membandingkan inode ketiganya. Setelah itu hapus
berkas asli dan jelaskan mengapa salah satu link masih berfungsi sementara yang lain tidak.
\end{challengebox}

% ------------------------------------------------------------------------
\section{Pencarian Berkas}
\label{sec:search}

Seiring waktu, sebuah sistem Linux bisa berisi ribuan berkas di ratusan direktori. Ketika
kamu perlu menemukan satu berkas konfigurasi atau mencari semua log yang berukuran besar,
menelusuri direktori satu per satu bukan pilihan yang masuk akal. Di sinilah perintah
pencarian berperan.

Perintah \cmd{find} melakukan pencarian secara langsung (\textit{real-time}) di filesystem.
Kamu bisa menentukan kriteria pencarian seperti nama berkas, tipe, ukuran, waktu modifikasi
terakhir, bahkan izin aksesnya. Setelah menemukan berkas yang cocok, \cmd{find} bisa
langsung menjalankan perintah tertentu pada setiap berkas yang ditemukan menggunakan opsi
\texttt{-exec}. Fleksibilitas ini menjadikan \cmd{find} salah satu perintah paling kuat
di Linux.

Perintah \cmd{locate} bekerja dengan cara yang berbeda. Ia mencari berdasarkan database
indeks yang dibangun sebelumnya oleh \cmd{updatedb}. Hasilnya sangat cepat karena tidak
perlu menelusuri filesystem secara langsung. Kelemahannya: jika berkas baru dibuat setelah
database terakhir diperbarui, \cmd{locate} tidak akan menemukannya. Gunakan \cmd{locate}
ketika kamu butuh jawaban cepat dan berkas yang dicari bukan yang baru saja dibuat.

Dua perintah tambahan yang berguna untuk mencari program: \cmd{which} menampilkan lokasi
lengkap sebuah program yang ada di \var{PATH}, dan \cmd{whereis} menampilkan lokasi binary,
halaman manual, dan kode sumber sebuah program.

\subsection*{Praktek 4.4 Cari Berkas dengan Kriteria Tertentu}

Tujuan praktek ini adalah menggunakan \cmd{find} dengan berbagai kriteria untuk menemukan
berkas secara efisien dan menggabungkannya dengan teknik redirection dari Bab~3.

\begin{enumerate}
    \item Buat beberapa berkas uji dengan nama dan ukuran yang berbeda:
\begin{lstlisting}[language=bash, caption={Membuat berkas uji untuk pencarian}]
cd ~/lab-os/chapter04-files
mkdir -p data-cari/logs data-cari/config data-cari/scripts
echo "isi log kecil" > data-cari/logs/app.log
echo "isi log besar" > data-cari/logs/error.log
yes "baris log panjang" | head -5000 > data-cari/logs/access.log
echo "host=localhost" > data-cari/config/app.conf
echo '#!/bin/bash' > data-cari/scripts/backup.sh
\end{lstlisting}

    \item Cari semua berkas \file{.log} di dalam direktori percobaan:
\begin{lstlisting}[language=bash, caption={Mencari berkas berdasarkan nama}]
find data-cari -type f -name "*.log"
\end{lstlisting}
    \textit{Amati:} Opsi \texttt{-type f} membatasi pencarian hanya pada berkas reguler, mengabaikan direktori. Opsi \texttt{-name} menerima pola dengan wildcard \texttt{*}.

    \item Cari berkas yang lebih besar dari 10 KB:
\begin{lstlisting}[language=bash, caption={Mencari berkas berdasarkan ukuran}]
find data-cari -type f -size +10k
\end{lstlisting}

    \item Cari semua berkas \file{.log} lebih besar dari 1 MB di \dir{/var/log/}
    dan simpan daftar hasilnya ke sebuah berkas menggunakan redirection:
\begin{lstlisting}[language=bash, caption={Mencari log besar dan menyimpan hasil}]
find /var/log -type f -name "*.log" -size +1M 2>/dev/null \
    > ~/lab-os/chapter04-files/log-besar.txt
cat ~/lab-os/chapter04-files/log-besar.txt
\end{lstlisting}
    \textit{Amati:} Pola \texttt{2>/dev/null} membuang pesan error ``Permission denied'' yang muncul ketika \cmd{find} mencoba masuk ke direktori yang tidak bisa diaksesnya. Teknik redirection ini sudah kamu pelajari di Bab~3.

    \item Gunakan \cmd{find} dengan \texttt{-exec} untuk menampilkan ukuran setiap berkas yang ditemukan:
\begin{lstlisting}[language=bash, caption={Menjalankan perintah pada setiap berkas yang ditemukan}]
find data-cari -type f -name "*.log" -exec ls -lh {} \;
\end{lstlisting}

    \item Gunakan \cmd{which} dan \cmd{whereis} untuk menemukan lokasi program:
\begin{lstlisting}[language=bash, caption={Menemukan lokasi program}]
which ls
which bash
whereis grep
\end{lstlisting}
\end{enumerate}

\begin{challengebox}
Gunakan \cmd{find} untuk mencari semua berkas \file{.log} di \dir{/var/log/} yang lebih
besar dari 1 MB. Arahkan hasilnya ke sebuah berkas bernama \file{laporan-log.txt}
di direktori \dir{~/lab-os/chapter04-files/}. Pastikan pesan error permission diabaikan
menggunakan \texttt{2>/dev/null}. Kemudian gunakan \cmd{grep} (dari Bab~2) dan \cmd{wc -l}
untuk menghitung berapa banyak berkas log besar yang ditemukan. Tuliskan jumlahnya sebagai
komentar di bawah perintahmu.
\end{challengebox}

% ------------------------------------------------------------------------
\section{Kompresi dan Pengarsipan}
\label{sec:archive}

Ketika kamu perlu mengirim puluhan berkas ke rekan tim, atau menyimpan cadangan sebuah
direktori sebelum melakukan perubahan besar, mengemas semua berkas menjadi satu paket
adalah langkah yang tepat. Proses ini disebut pengarsipan. Ketika paket tersebut juga
diperkecil ukurannya, proses itu disebut kompresi.

Di Linux, \cmd{tar} adalah perintah utama untuk pengarsipan. Nama \texttt{tar} berasal
dari \textit{tape archive}, karena dulu digunakan untuk menyimpan data ke pita magnetik.
Sekarang \cmd{tar} digunakan untuk menggabungkan banyak berkas dan direktori menjadi satu
berkas \texttt{.tar}. Secara default, \texttt{.tar} tidak terkompresi; ia hanya
menggabungkan. Untuk mengompresi sekaligus, kamu bisa menambahkan flag \texttt{-z}
agar \cmd{tar} menggunakan \cmd{gzip} secara otomatis, menghasilkan berkas \texttt{.tar.gz}.

Berikut flag \cmd{tar} yang paling sering digunakan:

\begin{itemize}
    \item \texttt{-c} : buat arsip baru (\textit{create})
    \item \texttt{-x} : ekstrak isi arsip (\textit{extract})
    \item \texttt{-t} : tampilkan isi arsip tanpa mengekstrak (\textit{list})
    \item \texttt{-f berkas} : tentukan nama berkas arsip (selalu dibutuhkan)
    \item \texttt{-v} : tampilkan nama berkas yang diproses (\textit{verbose})
    \item \texttt{-z} : gunakan kompresi gzip (untuk \texttt{.tar.gz})
    \item \texttt{-C direktori} : ekstrak ke direktori tertentu
\end{itemize}

Perintah \cmd{gzip} bekerja pada berkas tunggal: ia mengompresi satu berkas menjadi
\texttt{.gz} dan secara default menghapus berkas aslinya. Gunakan flag \texttt{-k} untuk
mempertahankan berkas asli, dan flag \texttt{-d} untuk dekompresi (setara dengan \cmd{gunzip}).

Perintah \cmd{zip} berguna ketika kamu perlu berbagi arsip dengan pengguna Windows atau
macOS, karena format \texttt{.zip} didukung secara native oleh semua sistem operasi populer.
Gunakan \cmd{zip -r} untuk mengarsipkan direktori secara rekursif, dan \cmd{unzip -l}
untuk melihat isi arsip tanpa mengekstraknya.

\subsection*{Praktek 4.5 Buat dan Ekstrak Arsip}

Tujuan praktek ini adalah membuat arsip terkompresi, memeriksa isinya, dan mengekstraknya
ke lokasi yang berbeda.

\begin{enumerate}
    \item Buat struktur direktori berisi beberapa berkas untuk diarsipkan:
\begin{lstlisting}[language=bash, caption={Menyiapkan direktori yang akan diarsipkan}]
cd ~/lab-os/chapter04-files
mkdir -p proyek/src proyek/config proyek/docs
echo "# Kode utama" > proyek/src/main.sh
echo "PORT=8080" > proyek/config/app.conf
echo "# Dokumentasi" > proyek/docs/readme.txt
ls -lR proyek/
\end{lstlisting}

    \item Buat arsip terkompresi dari direktori \dir{proyek/}:
\begin{lstlisting}[language=bash, caption={Membuat arsip .tar.gz}]
tar -czvf proyek-backup.tar.gz proyek/
ls -lah proyek-backup.tar.gz
\end{lstlisting}
    \textit{Amati:} Flag \texttt{-v} menampilkan setiap berkas yang ditambahkan ke arsip. Perhatikan ukuran arsip hasil kompresi dibandingkan ukuran direktori aslinya.

    \item Lihat isi arsip tanpa mengekstraknya:
\begin{lstlisting}[language=bash, caption={Melihat isi arsip tanpa ekstrak}]
tar -tzvf proyek-backup.tar.gz
\end{lstlisting}

    \item Ekstrak arsip ke direktori baru:
\begin{lstlisting}[language=bash, caption={Mengekstrak arsip ke direktori lain}]
mkdir -p hasil-ekstrak
tar -xzvf proyek-backup.tar.gz -C hasil-ekstrak/
ls -lR hasil-ekstrak/
\end{lstlisting}

    \item Kompres sebuah berkas tunggal menggunakan \cmd{gzip} sambil mempertahankan berkas aslinya:
\begin{lstlisting}[language=bash, caption={Kompresi berkas tunggal dengan gzip}]
yes "baris log" | head -10000 > catatan-panjang.txt
ls -lah catatan-panjang.txt
gzip -k catatan-panjang.txt
ls -lah catatan-panjang.txt catatan-panjang.txt.gz
\end{lstlisting}
    \textit{Amati:} Bandingkan ukuran berkas asli dengan versi yang sudah dikompres. Berapa persen penghematan ruang yang kamu dapatkan?

    \item Buat arsip ZIP untuk kompatibilitas lintas sistem operasi:
\begin{lstlisting}[language=bash, caption={Membuat arsip ZIP}]
zip -r proyek.zip proyek/
unzip -l proyek.zip | head -15
\end{lstlisting}
\end{enumerate}

\begin{challengebox}
Arsipkan seluruh direktori \dir{~/lab-os/chapter04-files/} menjadi sebuah berkas bernama
\file{chapter04-backup.tar.gz} yang disimpan di \dir{~/lab-os/}. Setelah selesai, gunakan
\cmd{tar -tzvf} untuk memverifikasi isi arsip dan pastikan semua subdirektori dan berkas
masuk ke dalam arsip. Kemudian gunakan pipa (\texttt{|}) untuk menghitung berapa banyak
berkas yang ada di dalam arsip tersebut menggunakan \cmd{wc -l}. Teknik pipa ini sudah
kamu pelajari di Bab~3.
\end{challengebox}

% ------------------------------------------------------------------------
\section{Rangkuman}
\label{sec:summary-filesystem}

Dalam bab ini, kamu telah mempelajari:

\begin{itemize}[noitemsep,topsep=2pt]
    \item \textbf{Hierarki FHS}: Linux mengatur semua berkas dalam satu pohon direktori dengan akar \dir{/}. Setiap direktori memiliki fungsi khusus: \dir{/etc/} untuk konfigurasi, \dir{/var/log/} untuk log, \dir{/home/} untuk data pengguna.
    \item \textbf{Jenis berkas}: Karakter pertama pada output \cmd{ls -l} menentukan jenis berkas. Tujuh jenis utama meliputi berkas reguler (\texttt{-}), direktori (\texttt{d}), tautan simbolis (\texttt{l}), dan berkas perangkat (\texttt{b}, \texttt{c}).
    \item \textbf{Links}: Hard link menunjuk ke inode yang sama dengan berkas asli, sedangkan symbolic link menunjuk ke path target. Perintah \cmd{ln} dan \cmd{ln -s} digunakan untuk membuat keduanya.
    \item \textbf{Pencarian berkas}: Perintah \cmd{find} mencari secara langsung di filesystem dengan berbagai kriteria seperti nama, tipe, ukuran, dan waktu modifikasi. Perintah \cmd{locate} mencari dengan cepat menggunakan database indeks yang perlu diperbarui secara berkala.
    \item \textbf{Kompresi dan pengarsipan}: Perintah \cmd{tar} menggabungkan banyak berkas menjadi satu arsip. Dengan flag \texttt{-z}, \cmd{tar} sekaligus mengompresi menggunakan \cmd{gzip}. Format \texttt{.zip} berguna untuk berbagi berkas lintas sistem operasi.
\end{itemize}

% ------------------------------------------------------------------------
\section{Latihan}
\label{sec:exercises-filesystem}

\textbf{Instruksi Umum}: Kerjakan seluruh latihan secara mandiri. Catat langkah penting, simpan tangkapan layar bila diperlukan, lalu rangkum hasilnya sebagai dokumentasi pribadi.

\subsubsection*{Latihan 4.1 Pemetaan Sistem File}
Jalankan \cmd{ls -lah /} dan pilih delapan direktori dari hasilnya. Untuk setiap direktori,
tuliskan fungsinya dalam satu kalimat dan berikan contoh satu berkas nyata yang ada di
dalamnya (verifikasi dengan \cmd{ls}). Sertakan bukti tangkapan layar untuk setiap
direktori yang kamu pilih.


\subsubsection*{Latihan 4.2 Hard Link dan Symbolic Link}
Buat satu berkas \file{sumber.txt} di \dir{~/lab-os/chapter04-tugas/}. Setelah itu:
\begin{enumerate}
    \item buat satu hard link bernama \file{sumber-hard.txt};
    \item buat satu symbolic link bernama \file{sumber-soft.txt};
    \item buktikan dengan \cmd{ls -li} mana yang berbagi inode;
    \item ubah isi berkas melalui hard link lalu tampilkan isi dari ketiga nama;
    \item hapus \file{sumber.txt} dan jelaskan keadaan kedua link yang tersisa.
\end{enumerate}
Sertakan semua perintah dan ringkasan pengamatanmu.


\subsubsection*{Latihan 4.3 Pencarian dan Pengarsipan}
Lakukan langkah-langkah berikut secara berurutan dan dokumentasikan hasilnya:
\begin{enumerate}
    \item Gunakan \cmd{find} untuk mencari semua berkas \file{.conf} di \dir{/etc/} yang
    berukuran kurang dari 5 KB. Arahkan hasilnya ke \file{config-kecil.txt}, buang
    error permission.
    \item Hitung berapa banyak berkas yang ditemukan menggunakan \cmd{wc -l}.
    \item Arsipkan seluruh direktori \dir{~/lab-os/chapter04-tugas/} menjadi
    \file{chapter04-tugas.tar.gz}.
    \item Verifikasi isi arsip dengan \cmd{tar -tzvf} dan hitung jumlah berkas di
    dalamnya menggunakan pipa ke \cmd{wc -l}.
\end{enumerate}


% ------------------------------------------------------------------------
\section{Referensi}
\label{sec:ref-filesystem}

\begin{itemize}
    \item Kerrisk, M. \textit{The Linux Programming Interface} (1st ed.). No Starch Press, 2010.
    \item Barrett, D., Silverman, R., Byrnes, R. \textit{Linux Pocket Guide} (3rd ed.). O'Reilly Media, 2016.
    \item Ward, B. \textit{How Linux Works} (3rd ed.). No Starch Press, 2021.
    \item Halaman manual: \texttt{man hier}, \texttt{man ls}, \texttt{man stat}, \texttt{man ln}, \texttt{man find}, \texttt{man tar}, \texttt{man gzip}.
\end{itemize}
